\ifdefined\gkver\else\def\gkver{student}\fi
\documentclass[\gkver]{gaokaozhenti}
\begin{document}

\chapter{2010年重庆理综}
%% paperType: 理综物理部分

\begin{enumerate}
\item
%% number: 14
%% paperName: 2010年重庆理综第 14 题 6 分
%% typeId: 1
%% type: 单选题
%% chapter: 机械波
%% point: 波的图象
%% method: 无
%% score: 6
%% degree: 850
%% duplicateId: 0
%% body:
一列简谐波在两时刻的波形如图中实线和虚线所示，由图可确定这列波的\xzanswer{C}

\onepicture[w=4.6cm]{figs/14.png}

\fourchoices[answer=C]
{周期}
{波速}
{波长}
{频率}

%% answer: C
%% memo:
\memoanswer{
由图象可确定波的波长。答案 C。
}

\item
%% number: 15
%% paperName: 2010年重庆理综第 15 题 6 分
%% typeId: 1
%% type: 单选题
%% chapter: 气体性质
%% point: 热力学第一定律
%% method: 无
%% score: 6
%% degree: 650
%% duplicateId: 0
%% body:
给旱区送水的消防车停于水平地面，在缓慢放水过程中，若车胎不漏气，胎内气体温度不变，不计分子间势能，则胎内气体\xzanswer{A}

\fourchoices[answer=A]
{从外界吸热}
{对外界做负功}
{分子平均动能减小}
{内能增加}

%% answer: A
%% memo:
\memoanswer{
温度是分子平均动能的标志，气体温度不变，分子平均动能不变，C 选项错误。不计分子势能，内能只由温度决定。气体温度不变，气体内能不变，D 选项错误。缓慢放水，胎内气膨胀对外做功，B 选项错误。由热力学第一定律知，从外界吸热，A 选项正确。
}

\item
%% number: 16
%% paperName: 2010年重庆理综第 16 题 6 分
%% typeId: 1
%% type: 单选题
%% chapter: 曲线运动
%% point: 双星问题
%% method: 无
%% score: 6
%% degree: 450
%% duplicateId: 0
%% body:
月球与地球质量之比约为 $1:80$，有研究者认为月球和地球可视为一个由两质点构成的双星系统，它们都围绕月地连线上某点 $O$ 做匀速圆周运动。据此观点，可知月球与地球绕 $O$ 点运动的线速度大小之比约为\xzanswer{C}

\fourchoices[answer=C]
{$1:6400$}
{$1:80$}
{$80:1$}
{$6400:1$}

%% answer: C
%% memo:
\memoanswer{
月球和地球构成的双星系统绕某点 $O$ 做匀速圆周运动，彼此间的万有引力提供向心力。由 $\frac{GMm}{l^{2}}=mv\omega$ 得 $\frac{v_{\text{月}}}{v_{\text{地}}}=\frac{M_{\text{地}}}{m_{\text{月}}}=\frac{80}{1}$，正确答案是 C。
}

\item
%% number: 17
%% paperName: 2010年重庆理综第 17 题 6 分
%% typeId: 1
%% type: 单选题
%% chapter: 交流电
%% point: 变压器
%% method: 无
%% score: 6
%% degree: 650
%% duplicateId: 0
%% body:
输入电压为 $220\UV$，输出电压为 $36\UV$ 的变压器副线圈烧坏，为获知此变压器原、副线圈匝数，某同学拆下烧坏的副线圈，用绝缘导线在铁芯上新绕了 5 匝线圈，如图所示，然后将原线圈接到 $220\UV$ 交流电源上，测得新绕线圈的端电压为 $1\UV$，按理想变压器分析，该变压器烧坏前的原、副线圈匝数分别为\xzanswer{B}

\onepicture[align=right, w=5cm]{figs/17.png}

\fourchoices[answer=B]
{$1100$，$360$}
{$1100$、$180$}
{$2200$、$180$}
{$220$、$360$}

%% answer: B
%% memo:
\memoanswer{
根据理想变压器 $\frac{u_{1}}{u_{2}}=\frac{n_{1}}{n_{2}}$，可求出变压器烧坏前的原、副线圈匝数。正确的答案是 B。
}

\item
%% number: 18
%% paperName: 2010年重庆理综第 18 题 6 分
%% typeId: 1
%% type: 单选题
%% chapter: 电路
%% point: 电容
%% method: 无
%% score: 6
%% degree: 450
%% duplicateId: 0
%% body:
某电容式话筒的原理示意图如图所示，$E$ 为电源，$R$ 为电阻，薄片 P 和 Q 为两金属极板，对着话筒说话时，P 振动而 Q 可视为不动，在 P、Q 间距离增大过程中，\xzanswer{D}

\onepicture[align=right, w=4cm]{figs/18.png}

\fourchoices[answer=D]
{P、Q 构成的电容器的电容增大}
{P 上电荷量保持不变}
{M 点的电势比 N 点的低}
{M 点的电势比 N 点的高}

%% answer: D
%% memo:
\memoanswer{
电容器 P、Q 间距增大，根据 $C=\frac{\varepsilon S}{4\pi kd}$，$C$ 减小，A 选项错误。电容器极板电压保持不变，根据 $C=\frac{Q}{u}$，P 上电荷量减少，电容器放电，方向为 M 到 N，B、C 选项错误，D 选项正确。
}

\item
%% number: 19
%% paperName: 2010年重庆理综第 19 题 6 分
%% typeId: 1
%% type: 单选题
%% chapter: 原子和原子核
%% point: 玻尔模型
%% method: 无
%% score: 6
%% degree: 650
%% duplicateId: 0
%% body:
氢原子部分能级的示意图如图所示，不同色光的光子能量如下所示：

\onepicture[align=right, w=3.2cm]{figs/19.png}

\begin{table}[htbp]
\centering
\begin{tblr}{|c|c|c|c|c|c|c|}
\hline
\SetCell[r=2]{c} 色光光子能量范围（$\UeV$） & 红 & 橙 & 黄 & 绿 & 蓝-靛 & 紫 \\
\hline
 & 1.61$\sim$2.00 & 2.00$\sim$2.07 & 2.07$\sim$2.14 & 2.14$\sim$2.53 & 2.53$\sim$2.76 & 2.76$\sim$3.10 \\
\hline
\end{tblr}
\end{table}

处于某激发态的氢原子，发射的光的谱线在可见光范围内仅有 2 条，其颜色分别为\xzanswer{A}

\fourchoices[answer=A]
{红、蓝-靛}
{黄、绿}
{红、紫}
{蓝-靛、紫}

%% answer: A
%% memo:
\memoanswer{
根据 $h\nu=E_{m}-E_{n}$，可知 A 选项正确。
}

\item
%% number: 20
%% paperName: 2010年重庆理综第 20 题 6 分
%% typeId: 1
%% type: 单选题
%% chapter: 光学
%% point: 全反射
%% method: 无
%% score: 6
%% degree: 450
%% duplicateId: 0
%% body:
如图所示，空气中有一折射率为 $\sqrt 2$ 的玻璃柱体，其横截面是圆心角为 $90^\circ$、半径为 $R$ 的扇形 OAB，一束平行光平行于横截面，以 $45^\circ$ 入射角照射到 OA 上，OB 不透光，若只考虑首次入射到圆弧 AB 上的光，则 AB 上有光透出部分的弧长为\xzanswer{B}

\onepicture[align=right, w=3.8cm]{figs/20.png}

\fourchoices[answer=B]
{$\frac{1}{6}\pi R$}
{$\frac{1}{4}\pi R$}
{$\frac{1}{3}\pi R$}
{$\frac{5}{12}\pi R$}

%% answer: B
%% memo:
\memoanswer{
光路图如图所示，由折射定律可求得设在 C 点恰好发生全反射，由 $\sin C=\frac{1}{n}$ 可求得 $C=45^\circ$，$\angle AOC=180^\circ-120^\circ-45^\circ=15^\circ$。弧 AB 上有光透出的部分弧长 CD 为 $\frac{90^\circ-30^\circ-15^\circ}{360^\circ}\cdot2\pi R=\frac{1}{4}\pi R$，正确的答案是 B。

\onepicture[w=3.8cm]{figs/20_an.jpg}
}

\item
%% number: 21
%% paperName: 2010年重庆理综第 21 题 6 分
%% typeId: 1
%% type: 单选题
%% chapter: 磁场
%% point: 带电粒子在磁场中的运动
%% method: 无
%% score: 6
%% degree: 450
%% duplicateId: 0
%% body:
如图所示，矩形 MNPQ 区域内有方向垂直于纸面的匀强磁场，有 5 个带电粒子从图中箭头所示位置垂直于磁场边界进入磁场，在纸面内做匀速圆周运动，运动轨迹为相应的圆弧，这些粒子的质量，电荷量以及速度大小如下表所示：

\begin{table}[htbp]
\centering
\begin{tblr}{|c|c|c|c|}
\hline
粒子编号 & 质量 & 电荷量（$q>0$） & 速度大小 \\
\hline
1 & $m$ & $2q$ & $v$ \\
\hline
2 & $2m$ & $2q$ & $2v$ \\
\hline
3 & $3m$ & $-3q$ & $3v$ \\
\hline
4 & $2m$ & $2q$ & $3v$ \\
\hline
5 & $2m$ & $-q$ & $v$ \\
\hline
\end{tblr}
\end{table}

由以上信息可知，从图中 a、b、c 处进入的粒子对应表中的编号分别为\xzanswer{D}

\onepicture[align=right, w=3.6cm]{figs/21.png}

\fourchoices[answer=D]
{3、5、4}
{4、2、5}
{5、3、2}
{2、4、5}

%% answer: D
%% memo:
\memoanswer{
根据 $qvB=m\frac{v^{2}}{r}$ 得轨迹半径 $r=\frac{mv}{Bq}$，结合表格中数据可求得 1--5 各组粒子的半径之比依次为 $0.5:2:3:3:2$，说明第一组正粒子的半径最小，该粒子从 MQ 边界进入磁场逆时针运动。由图 a、b 粒子进入磁场也是逆时针运动，则都为正电荷，而且 a、b 粒子的半径比为 $2:3$，则 a 一定是第 2 组粒子，b 是第 4 组粒子。c 顺时针运动，都为负电荷，半径与 a 相等是第 5 组粒子。正确答案 D。
}

\item
%% number: 22
%% paperName: 2010年重庆理综第 22 题 11 分
%% typeId: 4
%% type: 实验题
%% chapter: 电路
%% point: 伏安法测电阻
%% method: 无
%% score: 11
%% degree: 250
%% duplicateId: 0
%% body:
\textbf{（1）}（8 分）某同学用打点计时器测量做匀加速直线运动的物体的加速度，电源频率 $f=50\UHz$，在纸带上打出的点中，选出零点，每隔 4 个点取 1 个计数点，因保存不当，纸带被污染，如图所示，A、B、C、D 是依次排列的 4 个计数点，仅能读出其中 3 个计数点到零点的距离：$s_{A}=16.6\Umm$，$s_{B}=126.5\Umm$，$s_{D}=624.5\Umm$。

若无法再做实验，可由以上信息推知：

\begin{steps}
	\item
	相邻两计数点的时间间隔为\tkanswer[1.5cm]{0.1}$\Us$；
	\item
	打 C 点时物体的速度大小为\tkanswer[1.5cm]{2.5}$\Ums$（取 2 位有效数字）；
	\item
	物体的加速度大小为\tkanswer[3cm]{$\frac{(s_{D}-3s_{B}+2s_{A})f^{2}}{75}$}（用 $s_{A}$、$s_{B}$、$s_{C}$、$s_{D}$ 和 $f$ 表示）。
\end{steps}

\onepicture[w=14cm]{figs/22.png}

\textbf{（2）}（11 分）在探究小灯泡的伏安特性实验中，所用器材有：灯泡 L，量程恰当的电流表 A 和电压表 V，直流电源，滑动变阻器 $R$，电键 S 等，要求灯泡两端电压从 $0\UV$ 开始变化。

\begin{steps}
	\item
	实验中滑动变阻器应采用\tkanswer[1.5cm]{分压}接法（填“分压”或“限流”）；
	\item
	某同学已连接如图所示的电路，在连接最后一根导线的 C 端到直流电源正极之前，请指出其中仅有的 2 个不当之处，并说明如何改正：

	A．\tkanswer[5cm]{电键 $S$ 不应闭合，应处于断开状态}；

	B．\tkanswer[5cm]{滑动变阻器滑动触头 $P$ 位置不当，应将其置于 $b$ 端}。

	\onepicture[align=right, w=7cm]{figs/22a.png}
	\item
	电路连接正确后，分别测得两只灯泡 $L_{1}$ 和 $L_{2}$ 的伏安特性曲线如图中Ⅰ和Ⅱ所示。然后将灯泡 $L_{1}$、$L_{2}$ 与电池组（电动势和内阻均恒定）连成图所示电路。多次测量后得到通过 $L_{1}$ 和 $L_{2}$ 的电流平均值分别为 $0.30\UA$ 和 $0.60\UA$。

	A．在图中画出电池组路端电压 $U$ 和电流 $I$ 的关系曲线。

	B．由该曲线可知电池组的电动势为\tkanswer[1.5cm]{4.6}$\UV$，内阻为\tkanswer[1.5cm]{2.7}$\UO$。（取 2 位有效数字）
\end{steps}

\twopicture[align=right, showlabel=false, hA=4.4cm, hB=2.4cm]
{figs/22b.png}
{figs/22c.png}

%% answer: 
\jdanswer{
\textbf{（1）}0.1；2.5；$\frac{(s_{D}-3s_{B}+2s_{A})f^{2}}{75}$

\textbf{（2）}分压；电键 $S$ 不应闭合，应处于断开状态，滑动变阻器滑动触头 $P$ 位置不当，应将其置于 $b$ 端；$4.6\UV$，$2.7\UO$。
}
%% memo:
\memoanswer{
（1）①$T=0.02\times5=0.1\Us$

②$v_{c}=\frac{s_{D}-s_{B}}{2T}=\frac{(624.5-126.5)\times10^{-3}}{2\times0.1}\Ums=2.5\Ums$

③匀加速运动的位移特征是相邻的相等时间间隔内的位移以 $aT^{2}$ 均匀增大，有 $BC=AB+aT^{2}$，$CD=BC+aT^{2}=AB+2aT^{2}$，$BD=2AB+3aT^{2}$，所以

$a=\frac{(s_{D}-s_{B})-2\times(s_{B}-s_{A})}{3T^{2}}=\frac{(s_{D}-3s_{B}+2s_{A})f^{2}}{75}$

（2）①要求电压从 0 到额定电压变化，所以滑动变阻器必须采用分压接法。

②A 电键 S 不应闭合，应处于断开状态；B 滑动变阻器滑动触头 P 位置不当，应将其置于 b 端。

③描绘电源的伏安特性曲线要求外电阻变化测定对应的多组路端电压和电流，本实验中用两个小灯泡来改变外电阻获得两组电流值，然后在小灯泡的伏安特性曲线上查出对应的电压，用两个坐标点描绘图像。为描绘准确可以先进行理论计算，首先查出两坐标为（$0.30\UA$，$3.8\UV$）和（$0.60\UA$，$3.0\UV$），则内阻为 $r=\frac{3.8-3.0}{0.6-0.3}\UO=2.7\UO$，电动势为 $r=\frac{E-3.8}{0.3-0}$，$E=4.6\UV$，然后作出准确图像如图。

\onepicture[w=7cm]{figs/22b_an.png}
}

\item
%% number: 23
%% paperName: 2010年重庆理综第 23 题 16 分
%% typeId: 6
%% type: 计算题
%% chapter: 电磁感应
%% point: 电磁感应电路分析
%% method: 无
%% score: 16
%% degree: 450
%% duplicateId: 0
%% body:
法拉第曾提出一种利用河流发电的设想，并进行了实验研究。实验装置的示意图可用右图表示，两块面积均为 $S$ 的矩形金属板，平行、正对、竖直地全部浸在河水中，间距为 $d$。水流速度处处相同，大小为 $v$，方向水平。金属板与水流方向平行。地磁场磁感应强度的竖直分量为 $B$，水的电阻率为 $\rho$，水面上方有一阻值为 $R$ 的电阻通过绝缘导线和电键 K 连接到两金属板上。忽略边缘效应，求：
\begin{enumerate}
	\item
	该发电装置的电动势；
	\item
	通过电阻 $R$ 的电流强度；
	\item
	电阻 $R$ 消耗的电功率。
\end{enumerate}

\onepicture[align=right, w=6cm]{figs/23.png}

%% answer: 
\jdanswer{
\begin{enumerate}
	\item
	$Bdv$

	\item
	$\frac{BdvS}{\rho d+RS}$

	\item
	$\left(\frac{BdvS}{\rho d+RS}\right)^{2}R$

\end{enumerate}
}
%% memo:
\memoanswer{
（1）由法拉第电磁感应定律，有 $E=Bdv$。

（2）两板间河水的电阻 $r=\rho\frac{d}{S}$，由闭合电路欧姆定律，有 $I=\frac{E}{r+R}=\frac{BdvS}{\rho d+RS}$。

（3）由电功率公式，$P=I^{2}R$，得 $P=\frac{B^{2}d^{2}v^{2}S^{2}R}{(\rho d+RS)^{2}}$。
}

\item
%% number: 24
%% paperName: 2010年重庆理综第 24 题 18 分
%% typeId: 6
%% type: 计算题
%% chapter: 曲线运动
%% point: 平抛运动
%% method: 极值
%% score: 18
%% degree: 450
%% duplicateId: 0
%% body:
小明站在水平地面上，手握不可伸长的轻绳一端，绳的另一端系有质量为 $m$ 的小球，甩动手腕，使球在竖直平面内做圆周运动。当球某次运动到最低点时，绳突然断掉，球飞行水平距离 $d$ 后落地。如图所示。已知握绳的手离地面高度为 $d$，手与球之间的绳长为 $\frac{3}{4}d$，重力加速度为 $g$。忽略手的运动半径和空气阻力。
\begin{enumerate}
	\item
	求绳断时球的速度大小 $v_{1}$ 和球落地时的速度大小 $v_{2}$。
	\item
	问绳能承受的最大拉力多大？
	\item
	改变绳长，使球重复上述运动，若绳仍在球运动到最低点时断掉，要使球抛出的水平距离最大，绳长应是多少？最大水平距离为多少？
\end{enumerate}

\onepicture[align=right, w=6cm]{figs/24.png}

%% answer: 
\jdanswer{
\begin{enumerate}
	\item
	$v_{1}=\sqrt{2gd}$，$v_{2}=\sqrt{\frac{5}{2}gd}$

	\item
	$\frac{11}{3}mg$

	\item
	当 $l=\frac{1}{2}d$ 时，$x$ 有极大值，$x_{max}=\frac{2\sqrt{3}}{3}d$

\end{enumerate}
}
%% memo:
\memoanswer{
（1）设绳断后球飞行时间为 $t$，由平抛运动规律，有

竖直方向 $\frac{1}{4}d=\frac{1}{2}gt^{2}$，水平方向 $d=v_{1}t$

得 $v_{1}=\sqrt{2gd}$

由机械能守恒定律，有

$\frac{1}{2}mv_{2}^{2}=\frac{1}{2}mv_{1}^{2}+mg\left(d-\frac{3}{4}d\right)$

得 $v_{2}=\sqrt{\frac{5}{2}gd}$

（2）设绳子承受的最大拉力为 $T$，这也是球受到绳的最大拉力，球做圆周运动的半径 $R=\frac{3}{4}d$，

由向心力公式 $T-mg=m\frac{v_{1}^{2}}{R}$，得 $T=\frac{11}{3}mg$。

（3）设绳长为 $l$，绳断时球的速度大小为 $v_{3}$，绳承受的最大拉力不变，有 $T-mg=m\frac{v_{3}^{2}}{l}$，得 $v_{3}=\sqrt{\frac{8}{3}gl}$，

绳断后球做平抛运动，竖直位移为 $d-l$，水平位移为 $x$，时间为 $t_{1}$，有 $d-l=\frac{1}{2}gt_{1}^{2}$，$x=v_{3}t_{1}$，

得 $x=4\sqrt{\frac{l(d-l)}{3}}$，

当 $l=\frac{d}{2}$ 时，$x$ 有极大值 $x_{max}=\frac{2\sqrt{3}}{3}d$。
}

\item
%% number: 25
%% paperName: 2010年重庆理综第 25 题 19 分
%% typeId: 6
%% type: 计算题
%% chapter: 运动和力
%% point: 动量和能量
%% method: 无
%% score: 19
%% degree: 250
%% duplicateId: 0
%% body:
某兴趣小组用如图所示的装置进行实验研究。他们在水平桌面上固定一内径为 $d$ 的圆柱形玻璃杯，杯口上放置一直径为 $\frac{2}{3}d$，质量为 $m$ 的匀质薄圆板，板上放一质量为 $2m$ 的小物体。板中心、物块均在杯的轴线上，物块与板间动摩擦因数为 $\mu$，不计板与杯口之间的摩擦力，重力加速度为 $g$，不考虑板翻转。

\onepicture[align=right, w=6cm]{figs/25.png}

\begin{enumerate}
	\item
	对板施加指向圆心的水平外力 $F$，设物块与板间最大静摩擦力为 $f_{max}$，若物块能在板上滑动，求 $F$ 应满足的条件。
	\item
	如果对板施加的指向圆心的水平外力是作用时间极短的较大冲击力，冲量为 $I$，
	\begin{steps}
		\item
		$I$ 应满足什么条件才能使物块从板上掉下？
		\item
		物块从开始运动到掉下时的位移 $s$ 为多少？
		\item
		根据 $s$ 与 $I$ 的关系式说明要使 $s$ 更小，冲量应如何改变。
	\end{steps}
\end{enumerate}

%% answer: 
\jdanswer{
\begin{enumerate}
	\item
	$F>\frac{3}{2}f_{max}$

	\item
	$I>\frac{3}{2}m\sqrt{2\mu gd}$，$s=\frac{1}{2\mu g}\left(\frac{I-\sqrt{I^{2}-\frac{9}{2}\mu m^{2}gd}}{3m}\right)^{2}$，根据上式结果知：$I$ 越大，$s$ 越小。

\end{enumerate}
}
%% memo:
\memoanswer{
（1）设圆板与物块相对静止时，它们之间的静摩擦力为 $f$，共同加速度为 $a$，由牛顿运动定律，有：对物块 $f=2ma$，对圆板 $F-f=ma$，两物体相对静止，有 $f<f_{max}$，得 $F<\frac{3}{2}f_{max}$，相对滑动条件 $F>\frac{3}{2}f_{max}$。

（2）设冲击刚结束时圆板获得的速度大小为 $v_{0}$，物块掉下时，圆板和物块速度大小分别为 $v_{1}$ 和 $v_{2}$，由动量定理，有 $I=mv_{0}$，由动能定理，对圆板 $-2\mu mg\left(s+\frac{3}{4}d\right)=\frac{1}{2}mv_{1}^{2}-\frac{1}{2}mv_{0}^{2}$，对物块 $2\mu mgs=\frac{1}{2}(2m)v_{2}^{2}$，由动量守恒定律，有 $mv_{0}=mv_{1}+2mv_{2}$，要使物块落下，必须 $v_{1}>v_{2}$，解得：$I>\frac{3}{2}m\sqrt{2\mu gd}$，$s=\frac{1}{2\mu g}\left(\frac{I-\sqrt{I^{2}-\frac{9}{2}\mu m^{2}gd}}{3m}\right)^{2}$，分子有理化得 $s=\frac{1}{2\mu g}\left(\frac{\frac{3}{2}md}{I+\sqrt{I^{2}-\frac{9}{2}\mu m^{2}gd}}\right)^{2}$，根据上式结果知：$I$ 越大，$s$ 越小。
}

\end{enumerate}

\end{document}
